% Transcription — fonds Alexandre Grothendieck, shelfmark 154, batch 9 (pages 161-175).
% Established from the facsimile archives/batches/154/batch-09.pdf, cut from a
% copy of 154.pdf fetched through the project's own relay
% (https://grothendieck-relay.onrender.com/source/154.pdf); PDF page k =
% archivists' page 160+k, no cover sheet. Per the skill transcribe-grothendieck.
% The folder is 175 pages in nine batches, transcribed in parallel; this is
% the ninth and last (15 pages). Batches 1, 5 and 6 (transcripts/154/) were
% read first for their notation; notation otherwise follows folder 80
% (Systèmes de pseudo-droites).
%
% Pass: Opus 5.5 (claude-opus-5-5), 2026-09-26 — first pass, unchecked against the pages by a human.
%
% Dating and title. \dating is the folder's own date, copied verbatim from
% src/content/catalogue.ts; the folder belongs to the inventory group
% « Espaces stratifiés » (150-151), but since it has a date of its own the
% group's range is not added. Title from the catalogue, trailing full stop
% dropped; « [Système de pseudo-droites] » is bracketed there because the
% archivists supplied it.
%
% Paper. Every leaf is fan-fold computer listing paper written in landscape.
% Printed job banners and logs (a JCL header on 163; banners with 16 August
% 1982 on 164 and 166, 17 August 1982 on 172; an accounting log printed
% mirror-wise through the sheet on 168) are machine dates of the paper, not
% his, and are not transcribed.
%
% Pages skipped: none. Every leaf carries drawings of his.
% Pages summarised: none. Pages given only as a figure described in a French
% \note, with their few labels set: 161, 162, 163, 164, 166, 167, 168, 169,
% 172, 173, 174; figures with a short text or formula: 165 (a legend), 171
% (sequences of labels), 175 (a cycle of specialisations/generisations).
% Page transcribed from his hand: 170 (a list of axioms a)-h) for a pair of
% graphs K, L on a surface, beside two coloured drawings).
%
% His pagination and dates. No page carries a page number or a date of his.
% Numbered runs: on p. 170 the alineas a) b) (b) circled, written over a
% struck « b) Fam. ») c) (a first c) struck) d) e) f) g) h); on p. 175 the
% cycle 1 -> 2 -> 3 -> 4 -> 5 = 1; on pp. 171 and 173 circular sequences of
% labels 0-9, 0'-9' of pseudo-lines at the rim of a disc; on p. 162 rim
% labels 1-5, 1'-5'. All recorded as written.
%
% What the batch is about. Almost entirely coloured drawings, in the
% conventions of batches 5-6: discs with arrangements of pseudo-lines (red /
% orange), the moving line D (light green) sweeping them, supersingular bands
% and corridors (dark green), marked banks (blue), hatched triangles (yellow).
%   165    Legend: « Canal » -- place of the corners between marked (blue)
%          banks, and exceptional edges.
%   170    Axioms for (X, K, L): X a surface with boundary, K cellular, every
%          interior vertex of K of order 4 (boundary vertices of order 1);
%          an unfolding K' of K; a section of the sheets of the boundaries of
%          the immersion K' -> X; vertices of L of order 4; L ∩ K finite,
%          branches crossing; on each edge of K exactly one vertex of L; for
%          each interior vertex s of K, s not in L or s a vertex of L of
%          reduced order 2, eps_s = 0 or 1; for each face F of K,
%          sum_s (2 gamma_s + eps_s) + sum_a gamma_a = 2n, n independent of F.
%   171    Circular sequences of rim labels 3 2 1 0 4 5 ... 9' | 3 2 1 etc.,
%          i.e. the cyclic order of pseudo-lines on D~ before and after moves.
%   175    Cycle « 1 --spéc--> 2 --gén--> 3 --gén--> 4 --spéc--> 5 = 1 »,
%          drawn as successive positions of a line; a trefoil of three discs
%          with D_0, D_1 and alpha, beta, Delta.
%
% Notation as in batches 1, 5, 6: \mathcal{X} for his cursive X; X, K, L,
% K', \widetilde{K}' as he writes them; \gamma_s, \varepsilon_s, \gamma_a;
% S(K), A(K) for vertices and edges of K (his notation, read on p. 170).
%
% Readings to check first: p. 170, the whole of alinea b) (the struck « Fam. »
% and the line after it), the struck c), the end of c) « (connexe) ? » and
% the interlinear additions of d) and e); p. 165, the struck word of the
% second line and « arêtes » / « instables »; the sequences of p. 171.

\documentclass[11pt,a4paper]{article}
% Le préambule commun à toutes les transcriptions du fonds Grothendieck.
%
% Il tient en une idée : **l'incertitude se compose, elle ne se résout pas.**
% Une transcription qui lisse un mot douteux a détruit la seule chose pour
% laquelle on la faisait — un lecteur doit pouvoir savoir, à chaque endroit,
% ce qui était sur le feuillet et ce qui a été ajouté. Les six macros qui
% suivent sont donc l'appareil critique, et non de la décoration.
%
% Les mêmes commandes sont reconnues par `scripts/render.mjs`, qui produit la
% vue de lecture du site. Une macro ajoutée ici doit l'être aussi là-bas,
% faute de quoi le rendu échouera bruyamment — ce qui est voulu : un rendu
% silencieusement faux est pire qu'un rendu absent.

\ProvidesPackage{grothendieck}[2026/08/08 Transcriptions du fonds Grothendieck]

\usepackage{amsmath,amssymb,amsthm}
\usepackage{mathrsfs}
\usepackage{stmaryrd}
% Les diagrammes commutatifs sont partout dans ce fonds : c'est la langue
% même de la géométrie algébrique de Grothendieck, et une transcription qui
% les rendrait en prose ne transcrirait rien.
\usepackage{tikz-cd}
% Français par défaut — c'est la langue des feuillets — mais l'anglais est
% chargé aussi : la traduction et le résumé sont des documents anglais, et la
% typographie française (espaces avant les deux-points) y serait une faute.
% Ces documents basculent par \selectlanguage{english} en tête de corps.
\usepackage[english,french]{babel}
\usepackage{fontspec}
\usepackage{geometry}
\geometry{a4paper,margin=25mm,marginparwidth=28mm}
\usepackage{marginnote}
\usepackage{xcolor}
% ulem, not soul. `soul`'s \st and \ul are built on a character-by-character
% scan that XeLaTeX's font machinery breaks: the document fails with an
% undefined control sequence rather than degrading. ulem does the same two jobs
% and is Unicode-engine safe. `normalem` stops it hijacking \emph, which the
% transcriptions use for Grothendieck's own emphasis.
\usepackage[normalem]{ulem}
\usepackage{hyperref}
\hypersetup{colorlinks=true,linkcolor=black,urlcolor=black,pdfborder={0 0 0}}

% --- Métadonnées du lot -----------------------------------------------------
% Reprises telles quelles par le script de rendu, qui les affiche en tête de
% la vue de lecture. Elles fixent ce que le fichier prétend couvrir : sans
% elles, une transcription flotte sans point d'attache dans un fonds de seize
% mille feuillets.

\newcommand{\folder}[1]{\gdef\grothFolder{#1}}
\newcommand{\batch}[1]{\gdef\grothBatch{#1}}
\newcommand{\leaves}[2]{\gdef\grothFirst{#1}\gdef\grothLast{#2}}
\newcommand{\foldertitle}[1]{\gdef\grothFoldertitle{#1}}
\gdef\grothFolder{?}\gdef\grothBatch{?}\gdef\grothFirst{?}\gdef\grothLast{?}\gdef\grothFoldertitle{}

% --- Le résumé --------------------------------------------------------------
%
% Il ouvre la lecture modernisée, sous le titre « Résumé », et il est composé
% en petit. Ce n'est pas un ornement : c'est le seul passage du document qui
% ne soit pas des mathématiques, et un lecteur doit pouvoir le reconnaître
% avant de l'avoir lu — puis le sauter s'il connaît déjà le sujet, ou s'y
% arrêter s'il ne le connaît pas. Le filet qui le suit marque la reprise du
% corps.

\newenvironment{resume}
  {\small\setlength{\parskip}{0.45em}}
  {\par\medskip\hrule\medskip}

% --- Mots-clés ---------------------------------------------------------------
%
% En anglais, en clôture du résumé de la lecture modernisée : le vocabulaire
% moderne sous lequel chercher ce que les feuillets construisent. Le manifeste
% du site (`npm run manifest`) les extrait pour étiqueter la cote entière —
% cette ligne est la source unique des tags, il n'y a pas de fichier à côté.
\newcommand{\keywords}[1]{\par\smallskip\noindent
  {\footnotesize\textsc{Keywords} --- \textit{#1}}\par}

% --- L'appareil critique ----------------------------------------------------

% Le numéro de feuillet, en marge. C'est l'ancre qui permet de régler un
% désaccord sur une formule en regardant *un* feuillet plutôt qu'une chemise
% de six cents. Le site s'en sert aussi pour tourner les pages du fac-similé
% quand on fait défiler la transcription.
\newcommand{\leaf}[1]{%
  \marginnote{\footnotesize\color{gray!70}\textsf{#1}}%
  \label{leaf:#1}\ignorespaces}

% Illisible : on ne devine pas. Un mot inventé qui se lit comme les autres est
% le pire résultat possible.
\newcommand{\ill}{\textcolor{red!60!black}{[\,\ldots\,]}}

% Lecture douteuse : le mot est donné, et signalé comme incertain.
\newcommand{\uncertain}[1]{\uline{#1}}

% Ajout de l'éditeur — une lettre manquante, un mot sous-entendu.
\newcommand{\add}[1]{\textcolor{blue!50!black}{[#1]}}

% Biffé par Grothendieck lui-même. Ce qu'il a rayé fait partie du document :
% une rature dit souvent où la pensée a changé de direction.
\newcommand{\struck}[1]{\sout{#1}}

% Note du transcripteur, sur l'état matériel du feuillet.
\newcommand{\note}[1]{\par\smallskip{\small\color{gray!60!black}\textit{[#1]}}\par\smallskip}

% Note marginale de Grothendieck, distincte du corps du texte.
\newcommand{\marginal}[1]{\par\smallskip\noindent\hspace{1em}%
  {\small\color{gray!60!black}#1}\par\smallskip}

% --- Titre ------------------------------------------------------------------

\AtBeginDocument{%
  \begin{center}
    {\large\bfseries Fonds Alexandre Grothendieck --- cote n\textsuperscript{o}~\grothFolder}\\[2pt]
    {\itshape\grothFoldertitle}\\[4pt]
    {\small feuillets \grothFirst--\grothLast\ (lot \grothBatch)}\\[2pt]
    {\footnotesize\color{gray!60!black}Transcription établie d'après le fac-similé numérisé
     par l'Université de Montpellier.\\
     Les lectures incertaines sont soulignées, les illisibles notées [\,\ldots\,],
     les ajouts entre crochets.}
  \end{center}
  \vspace{1em}}

\endinput


\folder{154}
\batch{9}
\pages{161}{175}
\dating{1983-1984}
\watermark{Édition de démonstration}
\foldertitle{[Système de pseudo-droites] : notes manuscrites (1983-1984, s.d.)}

\begin{document}

\page{161}
\note{page de dessins sans texte suivi. En haut à gauche, au crayon, un disque traversé de trois pseudo-droites rouges qui se croisent près du centre, l'une dessinant une boucle étroite. En haut au milieu, un disque traversé de pseudo-droites rouges, plusieurs triangles coloriés en jaune, le bord doublé de tirets vert foncé et vert clair ; au bord, de petits numéros, en partie lisibles ($1$, $1'$, $2'$, $4'$, $5$). Au centre, un disque à bord rouge partagé par des arcs vert clair et vert foncé entre deux points $a$ (en bas) et $b$ (en haut), numéros $1$ à $5$ et $1'$ à $5'$ le long des arcs ; deux flèches en partent vers deux copies du même disque, les arcs doublés de pointillés et de tirets rouges, puis vers deux demi-disques, l'un fléché vers le haut le long des numéros $5'$, $4'$, $3'$, $2'$, l'autre vers le bas le long de $2$, $3$, $4$, $5$. En bas à gauche, au crayon, deux cercles superposés, l'arc supérieur $\alpha$ de $a$ à $b$, l'arc inférieur $\beta$ de $b$ à $a$, fléchés, avec deux orientations marquées $\omega$ ; à côté, un demi-disque dont le diamètre porte $\alpha$ au-dessus et $\beta$ au-dessous. En bas à droite, un rectangle bordé de rouge rempli d'arcs orange entrelacés, et un rectangle rouge coupé en deux par une bande vert clair}

\page{162}
\note{page de dessins sans texte. À gauche, deux disques accolés en huit, bordés de vert clair, traversés de pseudo-droites rouges, secteurs coloriés en jaune ; une double flèche noire relie un point du bord du disque supérieur à un point du disque inférieur. Au milieu, un disque dont une moitié du bord est vert foncé et l'autre vert clair, traversé de pseudo-droites rouges ; au bord, des numéros $0$, $1$, $2$, $3$, $4$, $5$ d'un côté et $0'$, $1'$, $3'$, $4'$, $5'$ de l'autre, les deux sommets opposés (marqués $0$, $1$ et $0$, $1'$) coupés d'un trait bleu. Au-dessous, un rectangle arrondi bordé de vert clair et de tirets rouges, traversé de pseudo-droites rouges, une ligne médiane vert foncé, deux traits bleus sur les côtés. À droite, au crayon, un cercle traversé d'un diamètre, avec au bord les marques $\alpha$ (deux fois), $1$, $1'$, $2'$, $3'$, $4'$, $5'$, $3$, $4$, $5$}

\page{163}
\note{un seul dessin en couleur : quatre courbes vert clair qui se coupent deux à deux, l'une formant une boucle ; les extrémités sont marquées d'un trait bleu, chaque point de croisement porte deux petites flèches noires, et chaque arc entre deux croisements est marqué d'un signe $+$ ou $-$}

\page{164}
\note{dessins en couleur. En haut à gauche, deux étoiles à quatre branches : deux droites vert clair qui se croisent, hérissées de petites droites orange, les bouts des branches fermés par des arcs orange et des coins bleus ou vert foncé ; en haut à droite, deux croix vert clair. En bas à gauche, deux ovales vert clair accolés en huit, doublés de tirets, traversés chacun de deux droites orange marquées $D$ et $D'$ ; au point de contact, une bande vert foncé et les points $x'$, $y'$ ; à côté, un ovale seul, les mêmes $D$, $D'$, $x'$, $y'$ ; au-dessous, un carré dont les coins sont $x_0$, $y_0$ (en haut, avec $x$, $y$ entre eux) et $x'_0$, $y'_0$ (en bas, avec $x'$, $y'$), traversé d'un croisement orange. Au milieu en bas, des faces vert clair hachurées reliées par des croisements orange, des arcs en tirets marqués $a$, $b$, et les points $u_0$, $u$, $y_0$, $y$. À droite, deux croquis : une droite vert clair de $x$ à $y$ sous un croisement orange, un arc brun $a$ qui l'enjambe, le segment $x_0 y_0$ noirci ; puis la même de $u$ à $v$, l'arc $b$, le segment $u_0 v_0$}

\page{165}
\note{grand dessin en couleur : un réseau de courbes vert clair aux extrémités marquées de traits bleus, coupées de petites croix orange (croisements de pseudo-droites) et de petits traits orange ; en travers, une courbe vert foncé longée d'un canal orange fait de losanges et d'échelons, qui l'accompagne sur toute sa longueur. En bas à gauche, la légende suivante, de sa main}

Canal ; : coins \struck{des} \struck{\ill{}}
\uncertain{lieu} des coins entre bords (bleus) marqués, et \uncertain{arêtes} exceptionnelles \uncertain{instables} (\uncertain{aux} \uncertain{sommets} \uncertain{stables}).
\note{un mot est écrit au-dessus du second mot biffé de la deuxième ligne, illisible ; la ponctuation « ; : » est telle qu'elle est lue}

\page{166}
\note{dessins en couleur sans texte : deux croquis de droites vert clair coupées de droites orange ; deux grands disques traversés de pseudo-droites rouges, secteurs coloriés en jaune, une partie du bord renforcée en vert foncé ou en noir, avec dans le premier une pseudo-droite en tirets noirs qui serpente et, au crayon, des hachures ; à droite, une bande horizontale faite de rectangles vert clair successifs, chacun traversé de droites rouges qui se croisent en losanges autour d'une ligne médiane noircie par segments}

\page{167}
\note{dessins en couleur sans texte suivi : plusieurs disques bordés de vert clair ou de vert foncé traversés de pseudo-droites rouges, certains secteurs coloriés en jaune, des sommets marqués d'un point ; des faisceaux de pseudo-droites rouges qui enjambent une droite vert clair en arc, ou se referment en fuseaux au-dessus et au-dessous d'un segment vert foncé ; dans l'un de ces fuseaux, les marques $D$, $D'$, $D''$ et une suite de petits $0$ et $1$ le long des droites ; une bande de fuseaux rouges séparés par des traits vert clair, avec de petits $0$ ; en bas à droite, en noir, une grille de droites dont une vert clair}

\page{168}
\note{à gauche, un grand disque bordé de vert clair traversé de pseudo-droites orange, de nombreux triangles et fuseaux coloriés en jaune ; au bord, des coins marqués de traits bleus et vert foncé, étiquetés $\alpha_0$, $\alpha_1$, $\alpha_2$, $\alpha_3$, $\alpha_4$, $\alpha_5$, $\alpha_6$, $\alpha_7$, chacun deux fois, et « $\alpha_8 = \alpha_0$ » ; un arc du bord entre $\alpha_2$ et $\alpha_3$ renforcé en vert foncé. À droite, une chaîne verticale de cercles au crayon enfilés sur deux pseudo-droites orange qui se tressent, avec des traits bleus aux points de passage et un gribouillis au crayon. En bas, un croisement de deux droites orange sur une droite vert clair, avec des flèches $a$, $b$ au-dessous et $a'$, $b'$ au-dessus, et une droite en tirets qui l'enjambe ; à gauche, au crayon, un fuseau étroit dans un cercle}

\page{169}
\note{deux grands disques en couleur sans texte, variantes l'un de l'autre : bord vert clair (dans le second, une moitié renforcée en vert foncé), pseudo-droites orange, triangles coloriés en jaune le long du bord, sommets au bord marqués d'un point noir ou d'un trait bleu, et des chemins en tirets noirs qui joignent certains coins en coupant les triangles}

\page{170}
\note{à gauche, deux grands dessins en couleur : un réseau de courbes vert clair aux extrémités bleues, coupées de croix orange, avec de petits arcs vert foncé ou bleus aux croisements ; au-dessous, un second réseau où les croisements portent des groupes de petits traits orange (un, deux, trois ou quatre), les faces numérotées $12$, $14$, $15$, $17$, $21$, $24$, $27$. Au bord gauche, un croquis : deux croix orange sur une droite vert clair, reliées par des tirets bleus ; en haut à droite, une droite vert clair coupée d'une croix orange, avec des arcs en tirets et des arcs vert foncé. Le texte occupe la colonne de droite et le coin inférieur gauche}

et $L$ (\uncertain{cellulaires})
\note{la ligne commence ainsi, en haut de la colonne ; le début de la phrase n'est pas sur la page}

$\mathcal{X}$, surface \uncertain{à} bord, $K$ cellulaire
\note{une ligne relie « $\mathcal{X}$ » à « surface »}

a) tout sommet \add{int.}\ de $K$ est d'ordre $4$ (sommet sur le bord d'ordre $1$)
\note{« int. » est écrit au-dessus de la ligne, avec un signe d'insertion}

\struck{b) Fam.}\ $K$ \ill{} \ill{} un \uncertain{déploiement} $K'$ de $K$, $K'$ \uncertain{contenant} \ill{} \ill{}
\note{le début de l'alinéa est surchargé : un « b) » et un mot souligné, lu « Fam. », sont barrés ; un petit mot est écrit, biffé, au-dessus de « $K'$ »}

(b) On a une \uncertain{section} des \ill{} \uncertain{dans} \uncertain{feuillets} $\widetilde{K}'$ de $K'$ des bords de l'immersion $K' \to X$
\note{le « b) » est entouré ; le tilde de $\widetilde{K}'$ est placé au-dessus du mot qui termine la ligne précédente, notre lecture le rapporte à $K'$}

\struck{c) \ill{} $L \subset X$, $1$-\uncertain{compl.}}

c) Tout sommet de $L$ est d'ordre \ill{} (\uncertain{connexe}) \ill{}
\note{un mot biffé au-dessus de « Tout » ; la fin de la ligne se perd dans le bord de la feuille, et « (\uncertain{connexe}) » suivi d'un signe est écrit à droite, au-dessus de l'alinéa d)}

d) $L \cap K$ fini, et \uncertain{à points} \ill{}
\note{deux mots l'un sur l'autre après « et », le mot inférieur souligné ; une accolade les rattache à une note à gauche}

\marginal{\ill{} $K \cap L$, chaque \uncertain{branche} de $K$ \uncertain{croise} chaque \uncertain{branche} de $L$}

e) Sur chaque arête de $K$, il y a exactement un sommet de $L$, \struck{\ill{}} \add{\ill{}}
\note{la fin de la ligne est biffée ; au-dessus, une addition interlinéaire, illisible, qui semble commencer par « sauf »}

f) Pour chaque sommet (\uncertain{int.}) $s$ de $K$, \uncertain{dans} les quatre arêtes qui en sortent, les nb.\ de pts de $L$ entre $s$ et le sommet de $K$ sur l'arête \uncertain{sont} \uncertain{indépendants} de l'arête \ill{}

g) Pour chaque sommet int.\ $s$ de $K$, ou bien $s \notin L$, ou bien $s$ est un sommet de $L$ d'ordre réduit $2$. Soit
\[
\varepsilon_s = \begin{cases} 0 & s \notin L \\ 1 & s \in L \end{cases}
\]
\note{alinéas g) et h) dans le coin inférieur gauche de la page ; « si » est écrit devant chaque condition ; devant $\varepsilon_s$, un signe surchargé, peut-être $\gamma_s$}

h) $\forall$ face $F$ de $K$, on a
\[
\sum_{s \in S(K) \cap F} (2\gamma_s + \varepsilon_s) + \sum_{a \in A(K),\ a \subset F} \gamma_a = 2n
\]
$n$ indép.\ de $F$
\note{sous le premier $\sum$, l'indice est surchargé ; « $S(K)$ », « $A(K)$ » et « $a \subset F$ » sont notre lecture}

\page{171}
\note{dessins en couleur et au crayon. En haut à gauche, une croix vert clair et un croisement orange sur une droite vert clair, une droite en tirets au-dessus. Au milieu, un grand cercle au crayon marqué de quatre arcs bleus et de flèches, et à l'intérieur un disque bordé de vert clair traversé de pseudo-droites noires, numérotées au bord $0$, $1$, $2$, $3$, …, $9$ d'un côté et $0'$, …, $9'$ de l'autre. En haut à droite, un grand disque traversé de pseudo-droites orange, triangles coloriés en jaune, avec au bord les numéros $1$, $2$, $3$, $1'$, $2'$, $3'$, $5$, $6'$, $8'$, $9$ et une flèche. En bas, un faisceau de droites orange et vert clair qui se tordent autour d'un croisement numéroté $1\,2\,3$ et $3\,2\,1$, avec une flèche ; un petit disque dont le bord est vert foncé sur une moitié et vert clair sur l'autre, numéroté $9$, $0$, $1$, $2$, $3$, $4$, … et $9'$, $8'$, …, $0'$, $1'$, $2'$, $3'$ ; enfin un nœud de pseudo-droites rouges serrées autour d'un segment noir hachuré, numérotées $1$, $2$, $3$. Les suites suivantes sont écrites entre les dessins}

$3\,2\,1\,0\,4\,5\,6\,7\,8\,9\,0'\,1'\,2'\,3'\,4'\,5'\,6'\,7'\,8'\,9' \mid 3\,2\,1$
\note{en haut, au milieu ; une accolade au-dessous va de $3$ à $3'$, un segment en est détaché ; le troisième chiffre est surchargé}

$1'\,2'\,3'\,0\,4'\,5'\,7'\,6'\,8'\,9'\,0'\,3\,2\,1\,4\,5$ \uncertain{$7\,6$} $8\,9 \mid 1'\,2'\,3'$
\note{une accolade sous le début de la suite ; « $7\,6$ » est surchargé}

$3\,2\,1 \mid 0\,4\,5\,6\,7\,8\,9\,0'\,1'\,2'\,3'\,4'\,5'\,6'\,7'\,8'\,9'$
\note{accolade sous $0 \dots 1'$}

$1'\,2'\,3'\,0\,4'\,5'\,7'\,6'\,8'\,9'\,0'\,3\,2\,1\,4\,5\,7\,6\,8\,9$
\note{accolade sous $3' \dots 3$}

$1\,0\,4\,5\,6\,7\,8\,9\,0'\,1'$
\note{au-dessus du petit disque du bas, avec une accolade}

\page{172}
\note{sur une feuille de listing imprimée, deux dessins en couleur sans texte : à gauche, trois disques accolés en chaîne, bordés de vert clair et de vert foncé, traversés de pseudo-droites rouges qui se croisent de l'un à l'autre, petits triangles coloriés en jaune ; à droite, une droite vert clair coupée de quatre droites rouges (deux croisées, deux simples), qui monte en passant par le deuxième croisement, doublée de tirets vert foncé et gris et de tirets bleus, avec une flèche vers la gauche et une vers la droite. Au crayon, un cercle en tirets et deux petits triangles}

\page{173}
\note{un grand dessin en couleur : un cercle vert clair entouré d'une couronne de pétales en tirets vert clair ; entre les pétales, des coins vert foncé et bleus sur le cercle ; dans la moitié supérieure, les pétales sont remplis de faisceaux de pseudo-droites rouges qui se tordent, avec des flèches de rotation et des numéros $0$, $1$, $2$, $3$, $4$, $5$, $6$, $7$ ; à l'intérieur du cercle, des traits numérotés $1$, $2$, $3$, $4$, $5$, $6$, $7$, $1'$, $2'$, $3'$, $5'$, $6'$, $7'$ et une flèche de rotation au centre. À droite, au crayon, un disque traversé d'arcs, et un disque où deux faisceaux se croisent, marqués $0$, $1$, $2$, $3$ en haut et $1'$, $2'$, $3'$ en bas ; à gauche, une petite étoile au crayon}

$1\,2\,3\,4\,5\,6\,7\,3\,2\,1\,4\,6\,5\,7$
\note{au crayon, sous le disque de droite ; lecture \uncertain{probable} des quatre derniers chiffres}

\page{174}
\note{dessins en couleur et au crayon sans texte : un disque rouge traversé de pseudo-droites rouges, bordé d'une couronne de pétales vert clair séparés par des coins vert foncé ; le même disque à moitié entouré, avec trois pétales au crayon, des coins bleus et vert foncé et, au bord, des numéros $1'$, $2'$, $3'$, $4'$, $5'$, $1$, $2$, $3$, $4$, $5$ ; un petit disque à trois pétales au crayon, un arc du bord fléché ; un disque traversé de trois droites rouges en triangle, un arc du bord bleu ; au crayon, une sphère coupée de grands cercles et deux polyèdres en perspective}

\page{175}
\[
1 \xrightarrow{\text{spéc}} 2 \xrightarrow{\text{gén}} 3 \xrightarrow{\text{gén}} 4 \xrightarrow{\text{spéc}} 5 = 1
\]
\note{en tête de la page ; la dernière étiquette est écrite « epéc », lue « spéc »}

\note{au-dessous, les positions successives de ce cycle : une droite rouge horizontale coupée d'un croisement de deux droites rouges et de deux droites verticales rouges ; une droite vert clair et des droites en tirets vert foncé, gris et bleu qui passent au-dessus ou au-dessous du croisement, leurs extrémités marquées $3$, $5 = 1$, $4$ et $2$, avec une flèche. À côté, deux cercles bleus autour de deux croisements, reliés par un arc. À droite, un disque traversé de pseudo-droites rouges, triangles coloriés en jaune, une bande en tirets rouges. En bas à gauche, un trèfle de trois disques bordés de vert clair, traversés de pseudo-droites rouges, secteurs coloriés en jaune, les bords marqués $D_0$, $D_1$ en alternance, et au centre les marques $\alpha$, $\beta$, $\Delta$, $\Delta_0$, $D_0$, $D_1$ ; au crayon, un disque sur lequel un petit segment vert porte deux cercles, et des ovales. Au milieu, un croquis de coins vert clair et vert foncé réunis par des traits. En bas à droite, une droite vert clair coupée de droites rouges, en partie en tirets, et d'une droite oblique rouge, deux segments de la droite verte noircis, des hachures au crayon}

\end{document}
